\section{Reference implementation}
\label{sec:impl}

\ampi{} is implemented in about 6{,}500 lines of dependency-free Python,
with 260 tests. The architecture borrows from MPICH and Open MPI, for the
reason those projects give: the layering is what lets the semantics outlive
any particular transport.

\subsection{Layering}

\begin{description}[leftmargin=1.2em,style=nextline]
\item[Device interface] The analogue of MPICH's Abstract Device
Interface~\cite{balaji2020mpich}. A device provides an at-least-once
message channel, a content-addressed blob store, a key-value store with
compare-and-swap, an append-only journal, and a clock. Notably it need
\emph{not} provide ordered delivery, connection state, or a broadcast
medium---the assumptions that make MPI devices hard to port. Everything
above the line is transport-independent, and the same test suite runs over
every device.

\item[Matching engine] Posted-receive and unexpected queues, MPI's
non-overtaking rule, deduplication, epoch filtering. Knows nothing about how
bytes move; this is the separation MPICH's repeated ADI refactors argue for.

\item[Communicators, collectives, RMA, I/O, fault tolerance] The protocol
proper, written against the matching engine.

\item[Profiler] Name-shifted interception of every operation, streaming to
the journal rather than buffering---if a rank dies, a buffered trace dies
with it, and the failure you most wanted to see is the one you lose.
\end{description}

Two devices are implemented. The \textbf{journal} device uses a shared
directory: message queues are directories, a send is one atomic
\code{rename}, a receive is a directory listing, and only compare-and-swap
needs exclusion. It requires nothing but a POSIX-ish filesystem, which is
deliberate---the ranks of an \ampi{} job are frequently processes we do not
control, each of which can run a shell command but cannot be handed a
socket. A directory is the lowest common denominator, as TCP was for early
MPI. The \textbf{inproc} device is the Nemesis analogue: identical
semantics, no syscalls, used for the test suite and for scaling studies
where filesystem latency would obscure the protocol's own costs.

\subsection{Consumption, not polling, is what commits}

One design point took us two attempts and is worth reporting, because it is
forced by an assumption most message-passing runtimes do not have to
question.

In a conventional runtime a rank is one process, so polling a message and
matching it are separated by microseconds and it is harmless to treat the
first as committing the second. An \ampi{} rank is frequently a
\emph{sequence} of short-lived processes---an agent that runs \code{ampi
recv}, thinks for two minutes, then runs \code{ampi send}---and every one of
them polls, matches at most one message, and exits. Our first implementation
marked messages consumed at poll time, so a process that polled six messages
and matched one silently destroyed the other five.

The fix is to separate the two: polling is speculative and in-process,
consumption is explicit and durable. That in turn requires per-rank
consumption state to survive across processes, and doing so naively---a set
of consumed message identifiers---would grow without bound. We keep instead
a per-\code{(context, source, destination)} watermark plus a small set of
out-of-order exceptions above it, compacted after every consumption, which
is $O(1)$ per peer in the common in-order case. The same durable state
carries the send sequence counters and the collective counter, which is what
lets a stateless endpoint sit in front of a stateful protocol.

\subsection{The command-line binding}

The binding that matters is not the Python one. \code{ampi} is a
command-line tool exposing every operation, so that

\begin{quote}
\emph{anything that can run a shell command can be an \ampi{} rank.}
\end{quote}

A rank's entire identity is two environment variables, \code{AMPI\_ROOT} and
\code{AMPI\_RANK}; everything else---the group, the rank table, the peers'
capabilities, the epoch---is discovered from the run directory at
\op{Init}, which is PMIx's design~\cite{castain2018pmix} scaled down.
Blocking operations block, which is what an agent wants: \code{ampi recv}
returns when the message arrives, and until then the agent is waiting at a
shell prompt.

\begin{figure}[t]
\begin{lstlisting}[language=ampishell]
# rank 3 of a 13-rank translation job
export AMPI_ROOT=/runs/book AMPI_RANK=3

ampi bcast   --root 0 --out spec.md
ampi scatter --root 0 --type json --out chunk.json
# ... read the chunk, propose terminology ...
ampi scan --exclusive --op ampi_first \
          --file proposal.json --type json --out prefix.json
# ... translate, adopting prefix.json ...
ampi gather  --root 0 --file result.json --type json
ampi progress
\end{lstlisting}
\caption{A complete rank program. Every line is a shell command, so the
executor can be a coding agent, a script, or a person. Errors are structured
JSON on stderr with a machine-readable \code{error} field, so an agent can
branch on \op{ERR\_TIMEOUT} versus \op{ERR\_REVOKED} without parsing prose.}
\label{fig:cli}
\end{figure}

This is the same bet MPI made on having bindings for the languages people
actually used. An agent protocol with only a Python binding is a Python
framework; a rank implemented by a coding agent, a shell script and a Python
program are indistinguishable on the wire and interoperate in one job.

\subsection{The launcher is out of scope}

\ampi{} specifies the wire-up and declines to specify the launch. A launch
plan is a run directory plus one prompt per rank; how those prompts become
running agents is the launcher's business. In our experiments the launcher
is Cursor's subagent API; it could be containers, serverless functions, or
people. \Cref{sec:discussion} reports what happened when the launcher turned
out to have a concurrency limit we had not accounted for.

\subsection{Observability}

Every operation emits enter/leave states and matched send/receive arrows,
following OTF2's and SLOG-2's event model, with token and currency counters
on every event. From that stream the tooling computes the standard MPI
views---a communication matrix, a Gantt timeline, per-operation time
breakdowns---plus two that are specific to this setting: cumulative context
pressure per rank over time, and a critical path measured in \emph{turns}
rather than seconds. The turn-based critical path is the more useful number,
because the ratio of total turns to critical-path turns is the harness's
achievable parallelism and does not depend on how fast the model happened to
be that day.
